\documentclass[12pt,a4paper]{article}

\usepackage{iftex}
\ifPDFTeX
  \usepackage[T1]{fontenc}
  \usepackage[utf8]{inputenc}
  \usepackage{lmodern}
\else
  \usepackage{fontspec}
  \setmainfont{Latin Modern Roman}
\fi
\usepackage[provide=*,serbian-latin]{babel}
\usepackage[a4paper,margin=2.5cm]{geometry}
\usepackage{microtype}
\usepackage{enumitem}
\usepackage{array}
\usepackage{graphicx}
\usepackage[hidelinks]{hyperref}
\IfFileExists{docs/diagrams/rendered/sr/component-architecture.png}{%
  \newcommand{\diagramassetpath}{docs/diagrams/rendered/sr/}%
}{%
  \newcommand{\diagramassetpath}{../docs/diagrams/rendered/sr/}%
}
\hypersetup{
  pdftitle={ASAP -- Projektovanje i implementacija AI MVP},
  pdfauthor={Mihajlo Madić 2119}
}

\setlength{\parindent}{0pt}
\setlength{\parskip}{0.7em}
\setlength{\emergencystretch}{2em}
\setlist{itemsep=0.4em,topsep=0.4em}

\newcommand{\placeholder}[1]{%
  \par\noindent\fbox{\parbox{0.95\linewidth}{\textit{Za dopunu: #1}}}\par
}

\title{Projektovanje i implementacija AI MVP\\[0.4em]
  \large Automatska Semantička Analiza Proizvoda (ASAP)}
\author{Mihajlo Madić, 2119\\
  MAS Računarstvo i informatika\\
  Veštačka inteligencija i mašinsko učenje}
\date{Inteligentni sistemi\\2026}

\begin{document}

\maketitle
\tableofcontents
\newpage

\section{Izveštaj 1: Predlog projekta}

\subsection{Osnovne informacije}

\begin{tabular}{>{\bfseries}p{0.28\linewidth}p{0.65\linewidth}}
Akronim & ASAP \\
Naziv rešenja & Automatska Semantička Analiza Proizvoda \\
Tim & Mihajlo Madić 2119 (MAS Veštačka inteligencija i mašinsko učenje) \\
\end{tabular}

\subsection{Kratak opis}

Mobilna aplikacija za automatsku identifikaciju i akviziciju podataka o proizvodima, koja koristi duboko učenje za prepoznavanje proizvoda u realnom vremenu i semantičke vektorske reprezentacije (engl. \textit{embeddings}) za generisanje personalizovanih preporuka i analitike zasnovane na semantičkoj sličnosti među proizvodima.

\subsection{Planirane AI tehnologije}

\begin{enumerate}
  \item \textbf{CNN (Computer Vision)} --- lokalizacija i dekodiranje barkodova proizvoda. Lagan model biće smešten na mobilnom uređaju i koristiće se za akviziciju vrednosti skeniranih barkodova.

  \item \textbf{Semantički embedding vektori (NLP)} --- generisanje vektorskih reprezentacija proizvoda u specijalizovanom servisu, na osnovu naziva, opisa, kategorije i drugih atributa prikupljenih preko eksternih API-ja.

  \item \textbf{Semantička pretraga} --- pronalaženje top-N najsličnijih proizvoda primenom metrika poput kosinusne sličnosti, uz opcionu MMR diversifikaciju radi izbegavanja redundantnih preporuka.

  \item \textbf{Personalizovane preporuke} --- formiranje vektorske reprezentacije korisnika kao centroida istorije skeniranih ili pregledanih proizvoda i generisanje preporuka na osnovu tog profila.

  \item \textbf{Personalna analitika (opciono)} --- klasterizacija proizvoda ili korisničkih interakcija i redukcija dimenzionalnosti pomoću PCA radi jednostavne vizuelne analitike u aplikaciji.
\end{enumerate}

Konkretna tehnologija lokalnog skenera barkoda još nije izabrana. Oznake CNN, TFLite i ML Kit iz početnih materijala predstavljaju kandidate koje treba proveriti u narednoj fazi, a ne potvrđenu implementacionu odluku.

\newpage

\subsection{Početna arhitektura}

\begin{figure}[ht]
\centering
\includegraphics[width=\textwidth,height=0.72\textheight,keepaspectratio]{\diagramassetpath component-architecture.png}
\caption{Prihvaćena planirana logička arhitektura komponenti. Otvorene tehnološke odluke su eksplicitno označene i dijagram ne predstavlja dokaz implementacije.}
\end{figure}

\clearpage
\subsection{Planirani tok od skeniranja do preporuke}

\begin{figure}[ht]
\centering
\includegraphics[width=\textwidth,height=0.78\textheight,keepaspectratio]{\diagramassetpath scan-to-recommendation-flow.png}
\caption{Planirani tok uspešnog zahteva i osnovnih ishoda kada barkod, metapodaci ili preporuke nisu dostupni.}
\end{figure}

\clearpage
\section{Izveštaj 2: Definisanje problema}

Ovaj deo se dopunjava tokom faze definisanja problema.

\subsection{Predlog rešenja}
\placeholder{promotivni opis rešenja do približno 150 reči}

\subsection{Ciljna grupa}
\placeholder{očekivani korisnici, inovativnost i razlika u odnosu na druga rešenja, do približno 300 reči}

\subsubsection{Korisnici}
\placeholder{primarne i sekundarne korisničke grupe}

\subsubsection{Glavne konkurentske prednosti}
\placeholder{prednosti u odnosu na postojeća rešenja}

\subsection{Detaljan opis rešenja}
\placeholder{koncept, osnovne funkcionalnosti i ograničenja}

\subsubsection{Osnovne funkcionalnosti}
\placeholder{funkcionalnosti nezavisne od AI komponenti}

\subsubsection{AI zasnovane funkcionalnosti}
\placeholder{funkcionalnosti zasnovane na računarskom vidu, embedding vektorima i preporukama}

\section{Izveštaj 3: Izbor tehnologija i skupova podataka}

Ovaj deo se dopunjava posle tehničkih eksperimenata i potvrde arhitekture.

\subsection{Tehnologije}
\placeholder{mobilna platforma, programski jezik, backend, skladište podataka i alati}

\subsection{AI tehnologije}
\placeholder{izabrani modeli, biblioteke, vektorska baza i metrike evaluacije}

\subsection{Skup podataka}
\placeholder{izvori podataka o proizvodima, licenca, obim, atributi i priprema podataka}

\section{Izveštaj 4: Razvoj MVP-a i dorade}

Prethodni delovi izveštaja dopunjavaju se u skladu sa odlukama i izmenama tokom razvoja.

\subsection{Odluke koje se odnose na prethodne faze}
\placeholder{promene zahteva, arhitekture, tehnologija ili skupa podataka i razlozi za njih}

\subsection{Prezentacija MVP-a nastavniku i drugim timovima}

\begin{tabular}{>{\bfseries}p{0.22\linewidth}p{0.7\linewidth}}
Datum & \textit{nije održana} \\
Komentari & \textit{za dopunu posle prezentacije} \\
\end{tabular}

\subsection{Korekcije i dorade}
\placeholder{povratne informacije, prioriteti i sprovedene izmene}

\subsection{Prezentacija konačne verzije MVP-a}

\begin{tabular}{>{\bfseries}p{0.22\linewidth}p{0.7\linewidth}}
Datum & \textit{nije održana} \\
Komentari & \textit{za dopunu posle prezentacije} \\
\end{tabular}

\section{Izveštaj 5: Testiranje sa potencijalnim korisnicima}

U ovoj fazi MVP se, u dogovoru sa nastavnikom, predstavlja potencijalnim korisnicima.

\begin{tabular}{>{\bfseries}p{0.28\linewidth}p{0.64\linewidth}}
Datum prezentacije & \textit{nije održana} \\
Podaci o korisniku & \textit{za dopunu} \\
Komentari & \textit{za dopunu} \\
\end{tabular}

\end{document}
